% ============================================================
% DOCUMENTCLASS
% ============================================================
\documentclass[12pt,oneside]{book}

% ============================================================
% METADATOS
% ============================================================
\newcommand{\universidad}{Universidad de Buenos Aires (UBA)}
\newcommand{\area}{Derecho Civil}
\newcommand{\materia}{Derechos Reales}
\newcommand{\titulo}{Derechos Reales de Disfrute sobre Cosa Ajena: Derecho Real de Habitación del Cónyuge y del Conviviente Supérstite}
\newcommand{\autor}{Isaac Edgar Camacho Ocampo}
\newcommand{\anio}{2025}

% ============================================================
% PAQUETES
% ============================================================
\usepackage[utf8]{inputenc}
\usepackage[T1]{fontenc}
\usepackage[spanish,es-nodecimaldot]{babel}

\usepackage[
    top=2.5cm,
    bottom=2.5cm,
    left=3cm,
    right=2.5cm,
    headheight=15pt
]{geometry}

\usepackage{amsmath}
\usepackage{amssymb}
\usepackage{mathtools}

\usepackage{graphicx}
\usepackage{float}
\usepackage{caption}
\usepackage{subcaption}

\usepackage{booktabs}
\usepackage{array}
\usepackage{longtable}
\usepackage{tabularx}

\usepackage{enumitem}

\usepackage{xcolor}
\usepackage[most]{tcolorbox}

\usepackage{fancyhdr}
\usepackage{ragged2e}

\usepackage[
    colorlinks=true,
    linkcolor=blue!50!black,
    citecolor=green!40!black,
    urlcolor=blue!60!black,
    pdftitle={\titulo},
    pdfauthor={\autor},
    pdfsubject={Apuntes universitarios},
    bookmarks=true
]{hyperref}

% ============================================================
% CONFIGURACIÓN
% ============================================================
\graphicspath{
    {./img/}
    {./graficos/}
    {./figuras/}
}

\setcounter{tocdepth}{2}

\setlength{\parindent}{0pt}
\setlength{\parskip}{0.5em}

\setlist[itemize]{
    topsep=0.3em,
    itemsep=0.2em
}

\setlist[enumerate]{
    topsep=0.3em,
    itemsep=0.2em
}

\clubpenalty=10000
\widowpenalty=10000

% ============================================================
% COMANDOS
% ============================================================
\newcommand{\abs}[1]{\left|#1\right|}
\newcommand{\norm}[1]{\left\lVert#1\right\rVert}
\newcommand{\vect}[1]{\mathbf{#1}}
\newcommand{\deriv}[2]{\frac{d#1}{d#2}}
\newcommand{\pderiv}[2]{\frac{\partial#1}{\partial#2}}

% ============================================================
% ENTORNOS / CAJAS
% ============================================================
\newtcolorbox{definition}[1]{
    enhanced,
    breakable,
    title={Definición: #1},
    fonttitle=\bfseries,
    colback=blue!3,
    colframe=blue!50!black,
    boxrule=0.8pt,
    arc=2mm
}

\newtcolorbox{important}{
    enhanced,
    breakable,
    title={Importante},
    fonttitle=\bfseries,
    colback=yellow!8,
    colframe=orange!70!black,
    boxrule=0.8pt,
    arc=2mm
}

\newtcolorbox{example}[1]{
    enhanced,
    breakable,
    title={Ejemplo: #1},
    fonttitle=\bfseries,
    colback=green!4,
    colframe=green!40!black,
    boxrule=0.8pt,
    arc=2mm
}

\newtcolorbox{formula}[1]{
    enhanced,
    breakable,
    title={Fórmula / Regla: #1},
    fonttitle=\bfseries,
    colback=purple!4,
    colframe=purple!50!black,
    boxrule=0.8pt,
    arc=2mm
}

\newtcolorbox{exercise}[1]{
    enhanced,
    breakable,
    title={Ejercicio: #1},
    fonttitle=\bfseries,
    colback=gray!5,
    colframe=gray!60!black,
    boxrule=0.8pt,
    arc=2mm
}

\newtcolorbox{note}{
    enhanced,
    breakable,
    title={Nota},
    fonttitle=\bfseries,
    colback=white,
    colframe=gray!50,
    boxrule=0.6pt,
    arc=2mm
}

% ============================================================
% CONFIGURACIÓN FINAL
% ============================================================
\pagestyle{fancy}
\fancyhf{}
\fancyhead[L]{\nouppercase{\leftmark}}
\fancyhead[R]{\materia}
\fancyfoot[C]{\thepage}
\renewcommand{\headrulewidth}{0.4pt}
\renewcommand{\footrulewidth}{0pt}

\fancypagestyle{plain}{
    \fancyhf{}
    \fancyfoot[C]{\thepage}
    \renewcommand{\headrulewidth}{0pt}
}

% ============================================================
% DOCUMENT
% ============================================================
\begin{document}

% ------------------------------------------------------------
% PORTADA
% ------------------------------------------------------------
\frontmatter

\begin{titlepage}

\thispagestyle{empty}

\begin{center}

\vspace*{1cm}

{\large \textsc{\universidad}}

\vspace{0.2cm}

{\large \area}

\vspace{3cm}

{\Huge \textbf{\materia}}

\vspace{1cm}

{\Large \titulo}

\vspace{0.8cm}

\textit{Estudiar con constancia es el camino para comprender con profundidad.}

\vspace{1.2cm}

\rule{0.70\textwidth}{0.5pt}

\vspace{0.5cm}

{\large Apuntes de estudio}

\vspace{0.5cm}

\rule{0.70\textwidth}{0.5pt}

\vfill

{\large \autor}

\vspace{0.3cm}

{\large \anio}

\end{center}

\vfill

\begin{flushleft}
\textit{Este es un modesto aporte a los estudiantes de ``Derechos Reales'' no pretende sustituir el material oficial de la materia.}
\end{flushleft}

\end{titlepage}

% ------------------------------------------------------------
% ÍNDICE
% ------------------------------------------------------------
\tableofcontents

\mainmatter

% ------------------------------------------------------------
% CONTENIDO
% ------------------------------------------------------------

% ============================================================
\chapter[Derechos reales de uso y de habitación]{Derechos reales de uso y de habitación}
% ============================================================

Este capítulo continúa el estudio de los \textbf{derechos reales de disfrute sobre cosa ajena}. Se analizan el \textbf{derecho real de uso} y el \textbf{derecho real de habitación}.

\section{Derecho real de uso}

\subsection{Definición y diferencia con el usufructo}

\begin{definition}{Derecho real de uso}
El derecho real de uso es el derecho de \textbf{usar y gozar} de una cosa que pertenece a otro en propiedad, pero \textbf{en forma limitada}.
\end{definition}

La diferencia fundamental con el \textbf{usufructo} reside en el límite. En el usufructo, la facultad de uso y goce es íntegra y comprende toda la facultad que tenía el titular dominial. % TODO: contenido incompleto o ambiguo en las notas originales (la redacción original de esta frase es confusa: «uso y goce como su fruto ario de toda la facultad que tenía el titular dominial»)
En cambio, en el derecho de uso aparece el \textbf{límite}: el titular puede usar y gozar de la cosa de otro, pero de manera limitada.

\subsection{El título constitutivo y el límite}

Los límites del derecho de uso deben estar \textbf{claramente consignados en el título constitutivo}. Este requisito es tan riguroso que, aun cuando se constituya un derecho real de uso sobre un inmueble mediante escritura pública y se lo denomine como tal, si el límite no consta claramente, \textbf{la figura se convierte en usufructo}.

\begin{quote}
\textbf{Artículo 2154 CCCN.} \textit{El uso es ese derecho real que consiste en usar y gozar de la cosa ajena, ya sea en una parte material o indivisa, en la extensión y con los límites establecidos en el título, que comparte con el usufructuario la carga de no alterar su sustancia.}
\end{quote}

La razón de la carga de no alterar la sustancia es que, una vez extinguido el derecho, la cosa debe restituirse en el estado en que se encontraba al momento de su constitución.

\subsection{Sujeto del derecho real de uso}

El derecho real de uso \textbf{solo puede constituirse a favor de una persona humana}. No necesariamente debe constituirse para cubrir las necesidades del usuario y su familia, como lo exigía la normativa anterior, pero sí debe estar ceñido a un límite al momento de su constitución. Ese límite se relaciona con la persona en particular, con el derecho que se le confiere y con el tiempo por el cual se le confiere.

\subsection{Régimen de los frutos}

Aun cuando el uso es limitado, si existieran únicamente frutos y, una vez abastecida la cantidad de frutos que corresponde al usuario y su familia, no hubiese más, \textbf{los frutos del usuario son preferidos a los del titular dominial}. Se trata, por lo tanto, de un derecho sumamente fuerte en tanto es un derecho real, aunque su ejercicio sea limitado.

\subsection{Transmisibilidad y derechos sobre la cosa}

El titular del derecho real de uso \textbf{no puede transmitirlo} ni constituir sobre la cosa derechos reales. El Código Civil y Comercial no aclara, sin embargo, lo relativo a los derechos personales.

\subsection{Inembargabilidad y ejecución por los acreedores}

Los derechos de disfrute son derechos patrimoniales e integran el patrimonio. No obstante, en el caso del derecho de uso, si está limitado a las necesidades del usuario y su familia, resulta \textbf{inembargable}.

\begin{quote}
\textbf{Artículo 2157 CCCN.} \textit{Si el uso se limita a las necesidades del usuario y su familia, el derecho es inembargable. El titular de este derecho no puede constituir sobre él derechos reales, ni puede trasmitirlo a título oneroso o gratuito.}
\end{quote}

\section{Derecho real de habitación}

\subsection{Definición y objeto}

\begin{definition}{Derecho real de habitación}
El objeto del derecho real de habitación es la \textbf{facultad de morar}, es decir, de habitar en un inmueble que pertenece a otro, ya sea en todo el inmueble o en una parte material. Solo puede constituirse a favor de una \textbf{persona humana}.
\end{definition}

\begin{note}
Queda planteada, para reflexión, la cuestión de la inscripción registral cuando el derecho recae no sobre todo el inmueble sino sobre una parte. % TODO: contenido incompleto o ambiguo en las notas originales (la duda se plantea pero no se resuelve)
\end{note}

\subsection{Norma aplicable}

\begin{quote}
\textbf{Artículo 2159 CCCN.} \textit{Supletoriamente se le aplican las normas del usufructo al derecho real de habitación.}
\end{quote}

\subsection{Limitaciones del derecho real de habitación}

\begin{table}[H]
\centering
\begin{tabularx}{\textwidth}{
    >{\raggedright\arraybackslash}p{0.30\textwidth}
    >{\raggedright\arraybackslash}X
}
\toprule
\textbf{Aspecto} & \textbf{Limitación} \\
\midrule
\textbf{Transmisión entre vivos} & No es transmisible por actos entre vivos. \\
\textbf{Transmisión mortis causa} & No es transmisible por causa de muerte. \\
\textbf{Constitución de derechos reales} & El habitador no puede constituir derechos reales sobre la cosa. \\
\textbf{Constitución de derechos personales} & El habitador no puede constituir derechos personales sobre la cosa. \\
\textbf{Ejecución por acreedores} & No es ejecutable por los acreedores del habitador. \\
\textbf{Gastos e impuestos} & El habitador debe abonar los impuestos, contribuciones y reparaciones necesarias para la conservación de la cosa, en proporción a la parte que ocupa. \\
\bottomrule
\end{tabularx}
\end{table}

\section{Cuadro Cornell}

\begin{longtable}{
    >{\raggedright\arraybackslash}p{0.25\textwidth}
    >{\raggedright\arraybackslash}p{0.63\textwidth}
}
\toprule
\textbf{Preguntas / palabras clave} & \textbf{Notas / respuestas} \\
\midrule
\endfirsthead
\toprule
\textbf{Preguntas / palabras clave} & \textbf{Notas / respuestas} \\
\midrule
\endhead
\bottomrule
\endfoot

\textbf{¿Qué es el derecho real de uso y en qué se diferencia del usufructo?} &
El derecho real de uso es el derecho de usar y gozar de la cosa ajena, pero en forma limitada. A diferencia del usufructo, que confiere una facultad íntegra de uso y goce, el derecho de uso exige que los límites estén claramente consignados en el título constitutivo. Si el título no establece el límite, la figura se convierte en usufructo (art.~2154 CCCN). \\[0.8em]

\textbf{¿Qué carga comparte el usuario con el usufructuario y por qué?} &
La carga de no alterar la sustancia de la cosa, porque al extinguirse el derecho la cosa debe restituirse en el estado en que se encontraba al constituirse. \\[0.8em]

\textbf{¿A favor de quién puede constituirse el derecho de uso?} &
Solo a favor de una persona humana. No es necesario que se constituya para las necesidades del usuario y su familia, pero sí debe estar ceñido a un límite relacionado con la persona, el derecho que se le confiere y el tiempo. \\[0.8em]

\textbf{¿Qué sucede con los frutos cuando el usuario tiene derecho a ellos?} &
Si existen únicamente frutos y, una vez satisfechas las necesidades del usuario y su familia, no sobran, los del usuario son preferidos a los del titular dominial. Es un derecho real fuerte aun siendo limitado. \\[0.8em]

\textbf{¿Es embargable el derecho real de uso?} &
No, cuando el uso está limitado a las necesidades del usuario y su familia, el derecho es inembargable (art.~2157 CCCN). Tampoco puede constituirse sobre él derechos reales ni transmitirse. \\[0.8em]

\textbf{¿Qué es el derecho real de habitación?} &
Es el derecho real que otorga la facultad de morar o habitar en un inmueble ajeno, en todo o en parte. Solo puede constituirse a favor de una persona humana. Supletoriamente se le aplican las normas del usufructo (art.~2159 CCCN). \\[0.8em]

\textbf{¿Qué limitaciones tiene el habitador?} &
No puede transmitir el derecho entre vivos ni por causa de muerte, no puede constituir derechos reales ni personales sobre la cosa, el derecho no es ejecutable por sus acreedores y debe pagar impuestos, contribuciones y reparaciones necesarias en proporción a la parte que ocupa. \\

\midrule

\multicolumn{2}{>{\raggedright\arraybackslash}p{0.89\textwidth}}{
\textbf{Resumen:} El derecho real de uso es un derecho de usar y gozar de cosa ajena de manera limitada; el límite debe constar en el título constitutivo, pues de lo contrario la figura se convierte en usufructo. Se constituye a favor de personas humanas, no es transmisible, y si se limita a las necesidades del usuario y su familia es inembargable. El derecho real de habitación otorga la facultad de morar en un inmueble ajeno, en todo o en parte, y está sujeto a limitaciones de transmisión, constitución de derechos y ejecución, además de la obligación de afrontar gastos en proporción a la parte ocupada.
} \\
\end{longtable}

% ============================================================
\chapter[Habitación del cónyuge supérstite]{Derecho real de habitación del cónyuge supérstite}
% ============================================================

\section{Naturaleza y ubicación en el Código}

El derecho real de habitación del cónyuge supérstite es un derecho real sobre cosa ajena, de los denominados de disfrute, con una particularidad singular: \textbf{no se encuentra regulado en el Libro Cuarto} del Código Civil y Comercial, donde se tratan conjuntamente los demás derechos reales, sino en el \textbf{Libro Quinto}, dedicado a la transmisión de los derechos por causa de muerte, es decir, las sucesiones. El artículo que lo regula es el \textbf{2383}.

Su ubicación en el Libro de las Sucesiones se explica porque este derecho se inserta en el conjunto de normas jurídicas que pretenden enfrentar las consecuencias de la finitud humana.

\begin{note}
El conocimiento de estos institutos resulta de aplicación frecuente en el ejercicio profesional: quien ejerce la abogacía, en el ámbito independiente, notarial, judicial o de asesoramiento, se encontrará con situaciones en las que un cónyuge o conviviente supérstite necesite ser informado y protegido en relación con la vivienda familiar. Se indica que el derecho de habitación era prácticamente desconocido incluso para los propios profesionales.
\end{note}

\section{Antecedente normativo}

\subsection{La necesidad de la protección}

\begin{example}{El caso de la tía Cota (año 1972)}
La tía Cota estaba casada con el tío Pedro, quien era viudo de su primer matrimonio y tenía un hijo ya adulto de esa unión anterior. Vivían en su casa propia, el tío trabajaba y no tenían grandes lujos ni necesidades. Hacia el año 1972 el tío Pedro falleció y se inició la sucesión. La casa en la que ambos habían vivido quedó dividida en un \textbf{50\% indiviso para la tía Cota y un 50\% indiviso para el hijo del tío Pedro}.

Cuando el hijo pidió la partición, el inmueble se puso en venta. Con el 50\% indiviso cobrado, descontados los gastos de sucesión, inmobiliaria y escritura, y con una pequeña pensión que apenas le alcanzaba para vivir, la tía Cota no pudo comprar nada y \textbf{quedó en la calle}.
\end{example}

Este modo de perder la vivienda, denominado ``pedir la parte'', sigue siendo una preocupación que viudos y viudas manifiestan con temor en la consulta profesional: qué sucederá si los otros herederos piden la parte.

\subsection{Artículo 3574 bis del Código Civil (Código de Vélez)}

Hasta el año 1974 no existía protección alguna para el cónyuge supérstite en situaciones como la descripta: cada vez que el hijo pedía la partición, el cónyuge viudo quedaba sin vivienda.

\begin{quote}
\textbf{Artículo 3574 bis CC (Ley 20.798, incorporado en 1974).} \textit{Incorporó el derecho real vitalicio y gratuito del cónyuge supérstite de habitar en el inmueble que fue asiento del hogar conyugal, siempre que fuera el único inmueble habitable del causante. Este derecho prevalecía aun ante el derecho de los herederos.}
\end{quote}

Esta norma presentaba limitaciones importantes:

\begin{itemize}
    \item El juez \textbf{no podía imponerlo de oficio}: solo se reconocía a solicitud de parte.
    \item Era un derecho \textbf{prácticamente desconocido}, tanto por sus titulares como por los propios colegas, por lo que, pese a existir desde 1974, viudas y viudos continuaban quedando sin vivienda por no invocarlo.
    \item El derecho \textbf{se extinguía} si el cónyuge supérstite contraía nuevas nupcias.
\end{itemize}

\section{Régimen actual: artículo 2383 del Código Civil y Comercial}

\begin{quote}
\textbf{Artículo 2383 CCCN.} \textit{El cónyuge supérstite tiene derecho real de habitación vitalicio y gratuito de pleno derecho sobre el inmueble de propiedad del causante que constituyó el último hogar conyugal y que a la apertura de la sucesión no se encontraba en condominio con otras personas. Este derecho es inoponible a los acreedores del causante.}
\end{quote}

Las características centrales del régimen se resumen en el siguiente cuadro:

\begin{table}[H]
\centering
\begin{tabularx}{\textwidth}{
    >{\raggedright\arraybackslash}p{0.24\textwidth}
    >{\raggedright\arraybackslash}X
    >{\raggedright\arraybackslash}p{0.16\textwidth}
}
\toprule
\textbf{Característica} & \textbf{Descripción} & \textbf{¿Novedad CCCN?} \\
\midrule
\textbf{Vitalicio} & Se ejerce durante toda la vida del cónyuge supérstite. & Sí \\
\textbf{Gratuito} & No se paga contraprestación por el uso del inmueble. & Sí \\
\textbf{De pleno derecho} & Se adquiere sin necesidad de solicitud judicial ni acto voluntario (art.~1894). & Sí (novedad) \\
\textbf{Sobre inmueble propio o ganancial} & No importa la calificación del bien: puede ser propio o ganancial del causante. & Sí (novedad) \\
\textbf{Sin requisito de único inmueble} & Ya no es necesario que sea el único inmueble habitable del causante. & Sí (novedad) \\
\textbf{No se extingue por nuevas nupcias} & No se pierde si el cónyuge supérstite se casa nuevamente o forma una unión convivencial. & Sí (novedad) \\
\textbf{Inoponible a acreedores del causante} & No puede hacerse valer frente a quienes el causante debía dinero. & Límite \\
\textbf{Requiere que no sea condominio} & El inmueble debe haber sido de propiedad exclusiva del causante al momento del fallecimiento. & Requisito \\
\bottomrule
\end{tabularx}
\end{table}

\section{Adquisición legal: artículo 1894 del Código Civil y Comercial}

\begin{quote}
\textbf{Artículo 1894 CCCN.} \textit{Se adquieren por imperio de la ley el derecho real de habitación del cónyuge supérstite y del conviviente supérstite, y los derechos reales sobre inmuebles que surgen de la ley.}
\end{quote}

Se trata de un caso de \textbf{adquisición legal} de los derechos reales: el derecho se adquiere por el solo hecho del fallecimiento del causante, sin necesidad de ningún acto jurídico adicional. Esto supera el principal problema práctico de la normativa anterior, que exigía la solicitud judicial.

\section{Tensión entre cónyuge supérstite y herederos}

Existe una tensión jurídica entre dos intereses. Por un lado, el cónyuge supérstite necesita conservar el único inmueble que constituye el asiento del hogar conyugal. Por el otro, los herederos tienen derecho a suceder a su causante, y sin duda lo tienen.

\begin{important}
El inmueble es uno: o se concede al cónyuge supérstite el derecho real de habitación de forma vitalicia, o se define a favor del derecho de los herederos. \textbf{El derecho ha optado por el derecho del cónyuge supérstite.}
\end{important}

\section{Constitucionalización del derecho privado}

El gran cambio de paradigmas producido con el Código Civil y Comercial de 2015 implica la denominada ``constitucionalización'' (o ``internalización'') del derecho privado: los principios y derechos reconocidos en la \textbf{Constitución Nacional} y en los \textbf{tratados de derechos humanos} penetran y modelan el derecho privado.

Bajo este paradigma, si se protege al cónyuge viudo, no se puede perder esa protección por el solo hecho de que se case nuevamente. Ello sería un castigo y una contradicción con la lógica protectoria del instituto. Por lo tanto, el derecho se defiende con independencia de si el cónyuge supérstite contrae nuevas nupcias, constituye una nueva unión convivencial o permanece solo.

\section{Cuadro Cornell}

\begin{longtable}{
    >{\raggedright\arraybackslash}p{0.25\textwidth}
    >{\raggedright\arraybackslash}p{0.63\textwidth}
}
\toprule
\textbf{Preguntas / palabras clave} & \textbf{Notas / respuestas} \\
\midrule
\endfirsthead
\toprule
\textbf{Preguntas / palabras clave} & \textbf{Notas / respuestas} \\
\midrule
\endhead
\bottomrule
\endfoot

\textbf{¿Qué es el derecho real de habitación del cónyuge supérstite?} &
Es un derecho real vitalicio y gratuito que se adquiere de pleno derecho (art.~1894 CCCN) sobre el inmueble que constituyó el último hogar conyugal y que era de propiedad exclusiva del causante al momento de la apertura de la sucesión (art.~2383 CCCN). Protege al cónyuge viudo frente al derecho de los herederos. \\[0.8em]

\textbf{¿Dónde está regulado y por qué allí?} &
En el Libro Quinto del CCCN (sucesiones), art.~2383, y no en el Libro Cuarto con los demás derechos reales. Se inserta en el conjunto de normas que enfrentan las consecuencias de la finitud humana. \\[0.8em]

\textbf{¿Por qué era necesaria la protección legal? ¿Cómo protegía la norma anterior?} &
Sin protección, cuando un heredero pedía la partición, el cónyuge supérstite podía quedar sin vivienda (caso de la tía Cota). El art.~3574 bis del Código de Vélez (1974) incorporó un derecho vitalicio y gratuito sobre el único inmueble habitable del causante que fue asiento del hogar conyugal, pero solo se reconocía a pedido de parte, era desconocido y se extinguía por nuevas nupcias. \\[0.8em]

\textbf{¿Cuáles son las novedades del CCCN respecto al régimen anterior (art.~3574 bis)?} &
1) Se adquiere de pleno derecho (antes se necesitaba solicitarlo). 2) No se extingue por nuevas nupcias ni por nueva unión convivencial. 3) No requiere que sea el único inmueble habitable del causante. 4) Puede recaer sobre bienes propios o gananciales del causante. \\[0.8em]

\textbf{¿Cómo se adquiere el derecho real de habitación del cónyuge supérstite?} &
Por adquisición legal, es decir, por imperio de la ley (art.~1894 CCCN). Se produce de pleno derecho con la apertura de la sucesión, sin necesidad de ningún acto judicial ni solicitud de parte. Esto supera el principal problema histórico: el desconocimiento del derecho por parte de los titulares y de los propios operadores jurídicos. \\[0.8em]

\textbf{¿A quién protege la ley en caso de conflicto entre cónyuge y herederos?} &
El inmueble es uno, y el derecho ha optado por el derecho del cónyuge supérstite, concediéndole la habitación vitalicia. \\[0.8em]

\textbf{¿Por qué se habla de ``constitucionalización del derecho privado'' en relación con este instituto?} &
El CCCN de 2015 incorporó al derecho privado los principios y derechos reconocidos en la Constitución Nacional y en los tratados de derechos humanos. Así, la protección del cónyuge supérstite no se pierde por contraer nuevas nupcias o formar una nueva convivencia, pues ello sería un castigo contrario a la lógica protectoria del instituto. \\

\midrule

\multicolumn{2}{>{\raggedright\arraybackslash}p{0.89\textwidth}}{
\textbf{Resumen:} El derecho real de habitación del cónyuge supérstite (art.~2383 CCCN) es vitalicio, gratuito y de pleno derecho, recae sobre el último hogar conyugal de propiedad exclusiva del causante y es inoponible a los acreedores del causante. Se adquiere por imperio de la ley (art.~1894 CCCN), sin necesidad de solicitud judicial, lo que supera las limitaciones del art.~3574 bis del Código de Vélez. Ante el conflicto con los herederos, el ordenamiento optó por proteger al cónyuge, en el marco de la constitucionalización del derecho privado.
} \\
\end{longtable}

% ============================================================
\chapter[Habitación del conviviente supérstite]{Derecho real de habitación del conviviente supérstite}
% ============================================================

\section{El reconocimiento de la unión convivencial en el CCCN}

La segunda gran novedad en esta materia es el reconocimiento de la \textbf{unión convivencial} como forma de familia protegida por el ordenamiento. Así como en algún momento se distinguía entre hijos legítimos e ilegítimos, distinción hoy superada, actualmente se reconoce que existen ``familias'': el término familia ya no está ligado exclusivamente a la institución matrimonial.

El Código reconoce la unión convivencial como una formación de familia, con o sin hijos, que puede durar muchos años y en la que uno de los convivientes puede quedar, de un día para el otro, sin donde vivir por el fallecimiento del otro.

\begin{quote}
\textbf{Artículo 527 CCCN (efectos de la convivencia).} \textit{El conviviente supérstite que carece de vivienda propia habitable o de bienes suficientes que aseguren el acceso a ésta, puede invocar el derecho real de habitación gratuito por un plazo máximo de dos años sobre el inmueble de propiedad del causante que constituyó el último hogar familiar y que a la apertura de la sucesión no se encontraba en condominio con otras personas. Este derecho es inoponible a los acreedores del causante. Se extingue si el conviviente supérstite constituye una nueva unión convivencial o contrae matrimonio.}
\end{quote}

\section{Diferencias con el régimen del cónyuge supérstite}

\begin{table}[H]
\centering
\begin{tabularx}{\textwidth}{
    >{\raggedright\arraybackslash}p{0.24\textwidth}
    >{\raggedright\arraybackslash}X
    >{\raggedright\arraybackslash}X
}
\toprule
\textbf{Aspecto} & \textbf{Cónyuge supérstite} & \textbf{Conviviente supérstite} \\
\midrule
\textbf{Duración} & Vitalicio. & Plazo máximo de 2 años. \\
\textbf{Extinción por nuevas nupcias o unión convivencial} & No se extingue. & Sí se extingue. \\
\textbf{Carencia de vivienda propia} & No lo exige. & Lo exige (condición de acceso o de extinción). \\
\textbf{Inmueble en condominio} & No opera si hay condominio. & No opera si hay condominio. \\
\textbf{Adquisición} & De pleno derecho (art.~1894). & De pleno derecho, pero con requisitos que los herederos pueden invocar. \\
\textbf{Inoponibilidad a acreedores} & Inoponible a acreedores del causante. & Inoponible a acreedores del causante. \\
\textbf{Regulación en el CCCN} & Libro Quinto, art.~2383. & Libro Segundo, art.~527. \\
\bottomrule
\end{tabularx}
\end{table}

\section{Debate sobre el artículo 527: adquisición de pleno derecho y requisitos}

El artículo 527 exige que el conviviente supérstite carezca de vivienda propia habitable o de bienes suficientes para proveerse una. Surge entonces la cuestión interpretativa de cómo se compatibiliza esta exigencia con la adquisición legal del derecho (art.~1894).

La posición expuesta en el material sostiene que el conviviente supérstite \textbf{adquiere el derecho de pleno derecho}: desde el momento del fallecimiento cuenta con los dos años que la ley le confiere para habitar gratuitamente el inmueble. Es facultad de los herederos invocar, durante ese período, que el conviviente tiene medios propios de vida o que posee otro inmueble en el cual vivir.

\begin{important}
El derecho del conviviente supérstite nace automáticamente con el fallecimiento, pero su subsistencia puede ser cuestionada por los herederos si acreditan que el conviviente posee recursos suficientes.
\end{important}

\section{La vivienda como derecho humano fundamental}

Todos estos institutos se encuadran en el reconocimiento del \textbf{derecho a la vivienda como derecho humano fundamental}. La ley ha tenido en cuenta este principio al regular el derecho real de habitación del cónyuge supérstite, de carácter vitalicio, y el del conviviente supérstite, limitado por un plazo de dos años, en correspondencia con la opción que cada familia realiza en el marco de su libertad: contraer matrimonio o constituir una unión convivencial.

\section{Cuadro Cornell}

\begin{longtable}{
    >{\raggedright\arraybackslash}p{0.25\textwidth}
    >{\raggedright\arraybackslash}p{0.63\textwidth}
}
\toprule
\textbf{Preguntas / palabras clave} & \textbf{Notas / respuestas} \\
\midrule
\endfirsthead
\toprule
\textbf{Preguntas / palabras clave} & \textbf{Notas / respuestas} \\
\midrule
\endhead
\bottomrule
\endfoot

\textbf{¿Qué reconocimiento introduce el CCCN respecto de las uniones convivenciales?} &
Reconoce la unión convivencial como forma de familia, con o sin hijos, que puede durar muchos años. El término familia ya no está ligado exclusivamente al matrimonio. \\[0.8em]

\textbf{¿Qué establece el artículo 527 CCCN sobre el conviviente supérstite?} &
Que el conviviente supérstite que carece de vivienda propia habitable o de bienes suficientes puede invocar el derecho real de habitación gratuito por un plazo máximo de dos años sobre el inmueble de propiedad del causante que fue el último hogar familiar y que no estaba en condominio. Es inoponible a los acreedores del causante y se extingue si constituye una nueva unión convivencial o contrae matrimonio. \\[0.8em]

\textbf{¿Qué diferencias existen entre el derecho del cónyuge supérstite y el del conviviente supérstite?} &
El del cónyuge es vitalicio y no se extingue por nuevas nupcias ni convivencia. El del conviviente está limitado a 2 años y sí se extingue si forma una nueva unión o contrae matrimonio. Además, el del conviviente exige que no tenga vivienda propia habitable ni bienes suficientes para acceder a una, y los herederos pueden invocar esa condición. \\[0.8em]

\textbf{¿El conviviente supérstite adquiere el derecho de pleno derecho?} &
Según la posición expuesta, sí: el derecho nace automáticamente con el fallecimiento y el conviviente tiene dos años para habitar gratuitamente el inmueble. Los herederos pueden invocar durante ese período que el conviviente tiene medios propios o que posee otro inmueble. \\[0.8em]

\textbf{¿Por qué se protege la vivienda de cónyuges y convivientes supérstites?} &
Porque la vivienda es reconocida como un derecho humano fundamental. La ley regula un derecho vitalicio para el cónyuge y uno limitado a dos años para el conviviente, en correspondencia con la opción que cada familia realiza: contraer matrimonio o constituir una unión convivencial. \\

\midrule

\multicolumn{2}{>{\raggedright\arraybackslash}p{0.89\textwidth}}{
\textbf{Resumen:} El CCCN reconoce la unión convivencial como forma de familia y, en el art.~527, otorga al conviviente supérstite un derecho real de habitación gratuito por un máximo de dos años, que se extingue si constituye una nueva unión o contrae matrimonio y exige carencia de vivienda propia habitable o de bienes suficientes. A diferencia del régimen del cónyuge (vitalicio y sin extinción por nuevas relaciones), el del conviviente es temporario. Ambos derechos responden a la protección de la vivienda como derecho humano fundamental frente a la muerte de uno de los integrantes de la familia, y su clave práctica es la adquisición legal de pleno derecho (art.~1894 CCCN).
} \\
\end{longtable}

\end{document}
